\begin{explanationbox}
\textbf{先看具体椭圆。} 取 $F_1(-3,0),F_2(3,0)$，固定距离和为 $10$。
\begin{center}\begin{tikzpicture}[font=\small,>=stealth]
\draw[black!30,rounded corners=2pt] (0,0) rectangle (14.4,8.2);
\node[anchor=west,font=\small\bfseries] at (.5,7.72) {同一椭圆上的两点：一段变长，另一段变短};
\begin{scope}[shift={(3.6,3.7)},scale=.55]
\draw[blue,very thick] (0,0) ellipse (5 and 4);
\draw[orange!85!black,very thick] (-3,0)--(0.000000000,4.000000000);\draw[violet,very thick] (3,0)--(0.000000000,4.000000000);
\fill (-3,0) circle (2pt) node[below] {$F_1$};
\fill (3,0) circle (2pt) node[below] {$F_2$};
\fill (0.000000000,4.000000000) circle (2pt) node[above] {$P$};
\end{scope}
\node at (3.6,1.05) {$PF_1+PF_2=5.00+5.00=10$};
\begin{scope}[shift={(10.7,3.7)},scale=.55]
\draw[blue,very thick] (0,0) ellipse (5 and 4);
\draw[orange!85!black,very thick] (-3,0)--(3.830222216,2.571150439);\draw[violet,very thick] (3,0)--(3.830222216,2.571150439);
\fill (-3,0) circle (2pt) node[below] {$F_1$};
\fill (3,0) circle (2pt) node[below] {$F_2$};
\fill (3.830222216,2.571150439) circle (2pt) node[above] {$P$};
\end{scope}
\node at (10.7,1.05) {$PF_1+PF_2=7.30+2.70=10$};

\node[anchor=west] at (.7,.55) {同尺度：$a=5,\ b=4,\ c=3$；右图长度取近似值，保留两位小数};
\end{tikzpicture}\end{center}

上方点的两段长都是 $5$；向右移动后，左段变长、右段变短，总长仍为 $10$。
这也是两钉加一根拉紧细绳画椭圆的原理：绳子两端固定在焦点，笔尖使两段绳的总长保持不变。
\textbf{只要求 $a>c$ 不够；动点还必须满足距离和等式。}
\end{explanationbox}
\clearpage
\begin{explanationbox}
\textbf{为什么 $a^2=b^2+c^2$？} 在垂直平分线上，两条焦半径相等，由距离和各得 $a$。
\begin{center}\begin{tikzpicture}[font=\small,>=stealth]
\draw[black!30,rounded corners=2pt] (0,0) rectangle (14.4,8.2);
\node[anchor=west,font=\small\bfseries] at (.5,7.72) {短轴端点：两条焦半径相等，所以每条都等于 $a$};
\begin{scope}[shift={(7.2,3.75)},scale=.58]
\draw[blue,very thick] (0,0) ellipse (5 and 4);
\draw[orange!85!black,very thick] (-3,0)--(0,4)--(3,0);
\draw[black!55] (-5,0)--(5,0);
\draw[violet,very thick] (0,0)--(0,4);
\draw[black!60] (0,.35)--(.35,.35)--(.35,0);
\node[left] at (0,2) {$b$};\node[below] at (1.5,0) {$c$};
\node[above right] at (1.5,2) {$a$};
\fill (0,4) circle (2pt) node[above] {$B$};
\fill (3,0) circle (2pt) node[below] {$F_2$};
\fill (-3,0) circle (2pt) node[below] {$F_1$};
\node[below left] at (0,0) {$O$};
\end{scope}

\node[anchor=west] at (.7,.55) {$BF_1=BF_2=a$；在直角三角形 $BOF_2$ 中，$a^2=b^2+c^2$};
\end{tikzpicture}\end{center}

于是直角三角形 $BOF_2$ 的斜边是 $a$，直角边是 $b,c$。
沿焦点连线取左右端点，则端点横坐标为 $\pm a$：两段长度是 $a-c$ 与 $a+c$。
所以长轴长是 $2a$，焦距是 $2c$，不能把两者当作同一个量。
\end{explanationbox}
\clearpage
